ModeBench makes a missing part of evaluation measurable: the range of distinct,
verified outcomes a model can produce for one problem. This matters because
different approaches can have different strengths. One may succeed where
another fails, or remain useful when constraints change. Understanding diversity
therefore gives a fuller picture of a model's capabilities than accuracy alone.

Our results show that replay can preserve alternatives while improving
held-out correctness. These controlled tasks provide a starting point for
studying diversity in broader mathematical reasoning. Replay retains only discovered,
stored modes; individual-mode survival remains unmeasured
(Appendix~\ref{app:ablations}). We have not shown that broader support causes
accuracy gains or helps solve later problems. The next step is to test which
alternatives help on richer mathematical tasks, and when.

